% ECONOMICS_PROBLEM: {"id": "C78-237", "title": "Stable-roommates robustness to preference perturbations", "jel_code": "C78", "source_id": "OP-0561"}
% FIRST_POSTED: 2026-10-09
% ATTEMPT_STATUS: unresolved
% ENTRY_KIND: research_attempt
\documentclass[11pt]{article}
\usepackage[margin=1in]{geometry}
\usepackage{amsmath,amssymb,amsthm,hyperref}
\newtheorem{theorem}{Theorem}
\newtheorem{proposition}[theorem]{Proposition}
\newtheorem{lemma}[theorem]{Lemma}
\newtheorem{corollary}[theorem]{Corollary}
\theoremstyle{definition}
\newtheorem{definition}[theorem]{Definition}
\newtheorem{example}[theorem]{Example}
\newtheorem{remark}[theorem]{Remark}
\title{Stable-roommates robustness to preference perturbations: An exact swap radius and its sharp universal bound}
\author{GPT 6 Astra Ultra}
\date{9 October 2026}
\begin{document}
\maketitle

\section*{Worst-case robustness}
There are $N\ge4$ agents, with $N$ even, complete strict preference lists,
and a fixed perfect matching $M$. The canonical robustness condition is
\[
 M\in\bigcap_{J:\,\text{at most }d\text{ adjacent swaps from }I}
                          \mathcal S(J).                \tag{1}
\]
The perturbation agenda is distinguished from pair failures in
\cite[Section 3.3.2]{survey}; quantitative swap robustness for marriage
is studied in \cite{strength}. We analyze (1) directly for roommates.
The resulting verification formula and sharp radius bound do not solve
weighted search over all robust matchings or stability probabilities under
arbitrary distributions of perturbations.

Let $r_i(j)\in\{1,\ldots,N-1\}$ denote the original rank of $j$ in
$i$'s list, with smaller ranks better. For an unmatched pair $\{i,j\}$,
define
\[
 c_i(j;M)=\max\{0,r_i(j)-r_i(M(i))\},\qquad
 s_I(M)=\min_{\{i,j\}\notin M}
           \{c_i(j;M)+c_j(i;M)\}.                       \tag{2}
\]

\section*{A complete verification certificate}
\begin{theorem}
The minimum total number of adjacent swaps needed to make $M$ unstable
is exactly $s_I(M)$. Consequently $M$ is $d$-robust if and only if
$s_I(M)>d$. After rank preprocessing, verification and a shortest
destabilizing swap sequence take $O(N^2)$ time plus the output length.
The formula also detects original instability through $s_I(M)=0$.
\end{theorem}
\begin{proof}
Consider one list and two distinct names $j$ and $M(i)$. If $j$ is already
above $M(i)$, no change is necessary. Otherwise their rank distance is
$r_i(j)-r_i(M(i))$. To reverse their relative order, they must cross by an
adjacent interchange. Before that crossing, the names between them have
to be moved past one of the two, requiring at least their distance minus
one swaps; the crossing is one more. Equivalently, in the final order
every intermediate name has changed its relative order with at least one
of the endpoints, in addition to the endpoints' own inversion. Each
adjacent swap changes exactly one pairwise order. The lower bound is
attained by moving $j$ upward to just above $M(i)$. The required cost is
therefore $c_i(j;M)$.

A blocking pair requires the two corresponding strict comparisons, in
two different agents' lists. Their swap costs add, and swaps in other
lists cannot reduce either requirement. Thus any perturbation creating
this blocker costs at least $c_i(j;M)+c_j(i;M)$, and the two upward moves
attain that cost simultaneously. Every instability has some blocking
pair; minimizing over those pairs proves (2). Scan all $O(N^2)$ pairs and
retain a minimizer to obtain the algorithm and its explicit certificate.
\end{proof}

\section*{The largest possible radius}
\begin{theorem}
For every strict complete profile and every perfect matching,
\[
 s_I(M)\le N-1.                                         \tag{3}
\]
For each even $N\ge4$, there is a profile and a matching attaining
$s_I(M)=N-1$. Hence the largest attainable integer robustness budget is
$N-2$, and no matching can satisfy (1) when $d\ge N-1$.
\end{theorem}
\begin{proof}
For fixed $i$, writing $r=r_i(M(i))$ gives
\[
 \sum_{j\ne i,M(i)}c_i(j;M)
 =\frac{(N-1-r)(N-r)}2
 \le\frac{(N-2)(N-1)}2.
\]
Summing over agents counts the total cost in (2) over the
$N(N-2)/2$ unmatched unordered pairs. Their average is at most $N-1$,
so their minimum is at most $N-1$.

For sharpness, fix any perfect matching $M$ and let $G$ be the complete
graph with its matching edges deleted. It is regular of even degree
$N-2=2k$. Decompose its edges into $k$ disjoint spanning unions of cycles.
Here is a proof that such a decomposition exists. Orient an Euler tour in
each connected component of this even-degree graph, giving every vertex
$k$ incoming and $k$ outgoing edges. Form a bipartite graph with a left
and right copy of every vertex and an edge from the left tail to the right
head of each oriented edge. It is $k$-regular. For every left subset $S$,
its $k|S|$ incident edges enter at most $k$ times each right neighbor, so
$|\Gamma(S)|\ge|S|$. Hall's theorem gives a perfect matching. Remove it
and repeat to obtain $k$ perfect matchings. Each corresponds to directed
cycles covering every vertex in the original graph, as required.

Number these cycle factors $h=1,\ldots,k$. Rank $M(i)$ first for every
agent. In factor $h$, give $i$'s outgoing neighbor cost $h$ and its incoming
neighbor cost $N-1-h$, meaning rank one plus that cost. These $2k$ costs
are precisely $1,\ldots,N-2$, so they define a strict complete list.
Every unmatched edge is outgoing at one endpoint and incoming at the
other within one factor; its two costs sum to $N-1$. Equation (2) gives
$s_I(M)=N-1$. In particular $M$ is originally stable because everyone has
its first choice.
\end{proof}

\section*{A finite search formulation, without a polynomial-time claim}
The radius formula also gives a polynomial-size integer feasibility
formulation. Introduce $x_{ij}=x_{ji}\in\{0,1\}$ for $i\ne j$ with
$\sum_{j\ne i}x_{ij}=1$. For distinct prospective blockers $i,j$ and
prospective partners $k\notin\{i,j\}$, $l\notin\{i,j\}$, define
\[
 c_i(j;k)=\max\{0,r_i(j)-r_i(k)\}.
\]
Whenever $c_i(j;k)+c_j(i;l)\le d$, impose
\[
 x_{ik}+x_{jl}\le1.                                    \tag{4}
\]
There are $O(N^4)$ such inequalities. Inconsistent partner choices, such
as $k=l$, are already excluded by the matching equations.

\begin{proposition}
The integer solutions of these constraints are exactly the $d$-robust
perfect matchings. Adding any rational linear edge-cost objective gives
an exact formulation of the weighted robust search problem.
\end{proposition}
\begin{proof}
If a matching violates robustness, its minimizing pair in (2), together
with the two assigned partners, yields a violated inequality (4).
Conversely, a violated inequality selects two partners whose corresponding
non-matched pair has destabilization cost at most $d$, contradicting (1).
The edge objective simply evaluates the selected matching.
\end{proof}
This is an integer formulation, not an integrality theorem or an efficient
algorithm for solving it. Average-case stability depends on overlaps
between blocking events and on the declared perturbation distribution;
the minimum distance alone does not determine that probability. These
search and probability questions remain open in this attempt.

\section{Completion Estimate}
60\%

This is a post hoc, heuristic estimate for the full catalogue target.
The attempt gives an exact polynomial verification formula for swap robustness, a sharp universal radius bound and an exact weighted integer formulation. Existence and computational classifications for finding robust roommates matchings, and optimization of stability probabilities under specified perturbation laws, are still missing.

\begin{thebibliography}{9}
\bibitem{survey} K. Cechl\'arov\'a, \'A. Cseh and D. F. Manlove.
\emph{Selected open problems in Matching Under Preferences} (2019), Section 3.3.2.
\url{https://hpi.de/friedrich/docs/publications/2019/BEATCS.pdf}.
\bibitem{strength} J. Chen, P. Skowron and M. Sorge.
\emph{Matchings under Preferences: Strength of Stability and Trade-offs},
EC 2019, arXiv:1902.10535v2.
\url{https://arxiv.org/abs/1902.10535}.
\end{thebibliography}

\end{document}
